\section{Erd\H{o}s Problem 175: square factors of central binomial coefficients}

\subsection{1) Statement and scope}
Show that $\binom{2n}{n}$ is not squarefree for every $n\geq5$.
The assertion has been proved by Granville--Ramar\'e and independently
by Velammal. This attempt proves it for every $n$ which is not a power
of $2$, and also for an infinite family of powers of $2$. It does not
reconstruct the treatment of every remaining power of $2$.

\subsection{2) The binary-digit reduction}
For a prime $p$, count the multiples of $p,p^2,\ldots$ in a factorial:
\[
 v_p(n!)=\sum_{j\geq1}\left\lfloor\frac n{p^j}\right\rfloor.
\]
If $n=\sum_i\epsilon_i2^i$ with $\epsilon_i\in\{0,1\}$, exchanging
the finite sums gives $v_2(n!)=n-\sum_i\epsilon_i=n-s_2(n)$.
Because $s_2(2n)=s_2(n)$,
\[
 v_2\binom{2n}{n}
 =(2n-s_2(2n))-2(n-s_2(n))=s_2(n).
\]
Every positive integer other than a power of $2$ has at least two binary
ones. Thus $4$ divides the binomial coefficient for every such $n$.
This covers all but $O(\log X)$ values of $n\leq X$, but that density
statement does not settle the universal quantifier.

\subsection{3) A further infinite family of exceptional inputs}
The factorial formula also yields
\[
 v_3\binom{2n}{n}
 =\sum_{j\geq1}
 \left(\left\lfloor\frac{2n}{3^j}\right\rfloor
             -2\left\lfloor\frac n{3^j}\right\rfloor\right).
\]
Each summand is either zero or one, and it equals one exactly when
the residue of $n$ modulo $3^j$ is at least $(3^j+1)/2$.

Set $n=2^k$ with $k\equiv3$ or $5\pmod6$. Such $k$ are odd, so
$n\equiv2\pmod3$ and the first summand is one. The powers of $2$
modulo $9$ have period $6$, with residues $2,4,8,7,5,1$.
For the two specified residue classes of $k$, the residue is respectively
$8$ or $5$, so the second summand is also one. Hence
\[
 9\mid\binom{2^{k+1}}{2^k}\qquad(k\equiv3,5\pmod6).
\]
All later summands are nonnegative. This argument therefore supplies
genuine square divisors for these infinitely many powers of $2$.

\subsection{4) Gap and outcome}
The remaining inputs in this proof are powers $n=2^k$, $k\geq3$, with
$k\equiv0,1,2,4\pmod6$. Some of them have two carries at later powers
of $3$, or square factors at other primes, but no uniform argument for
them has been established here. Finite computation cannot replace that
step. The stronger questions about the largest prime-power divisor,
mentioned in the problem's commentary, are also not proved here.

\textbf{Outcome of this attempt: UNRESOLVED in full.} The exact valuation
formula and the two families above are proved; the original all-$n$
theorem is known from the literature. No new resolution is claimed.
\begin{itemize}
\item A. Granville and O. Ramar\'e, \emph{Explicit bounds on exponential
sums and the scarcity of squarefree binomial coefficients}, Mathematika
43 (1996), 73--107; bibliographic links are on the problem page.
\item T. F. Bloom, Problem 175,
\texttt{https://www.erdosproblems.com/175}, checked 23 September 2026.
\end{itemize}
The percentage is a subjective proof-progress estimate, not the density
of inputs covered; the exceptional powers are the main difficulty.
\subsection{5) COMPLETION ESTIMATE}
\textbf{COMPLETION: 25\%}
